\chapter{Rate Theory In More Detail}

This chapter is organized into five sections. The first involves a revision of the
traditional concepts of dormancy and particularly "breaking dormancy." The other four
involve different aspects of rates of seed germination.

\section{Dormancy and "Breaking Dromancy" In Seeds and Plants}

It has been known for many years that deciduous trees and shrubs as well as
corms and bulbs require a period of time at temperatures around 32-45 during winter
in order for stem and leaf growth to resume in spring. The same chilling period is
needed for the flowering of deciduous trees and shrubs in spring. Times and
temperatures for such chilling periods have been determined in detail for important
commercial crops such as cherries and other fruit trees. As already noted in Chapter
5, these chilling requirements are analogous to the chilling requirements for the
germination of certain seeds.

\textit{A cardinal principle of physical organic chemistry is that if condition A is
required in order for a process to proceed under conditions B, critical chemistry must
have been taking place in conditions A}. Applying this principle to the chilling
requirements described above, it means that during the chilling period critical
chemical reactions were occurring. These chemical reactions are the destruction of
growth blocking systems whether this be the destruction of a specific inhibitor or
something more subtle such as the restructuring of a membrane or a protein. In either
situation it is a chemical change that is taking place.

Not only must these chemical changes be taking place, but they are taking
place at their \textbf{maximum metabolic rate}. Clearly they are taking place at
temperatures around 40, and their rates must be in effect zero at temperatures around
70. This is again an example of the enormous rate changes found in certain biological
processes, and this phenomenon in terms of rate theory is discussed in the fourth
section of this Chapter.

The traditional view that trees and shrubs are dormant in late fall, winter, and
early spring is at best misleading and completely wrong in terms of the destruction of
growth inhibitors. In terms of the latter the trees and shrubs are at their \textbf{maximum
metabolic activity} in late fall, winter, and spring. It is of course true that in respect to
vegetative growth the trees and shrubs are dormant, and in this sense they are
dormant. The problem is that the dormancy in respect to vegetative growth has misled
bioligists into thinking that all metabolic activity has ceased or at least fallen to a low
level. It cannot be emphasized enough that it is not just a matter of the blocking
systems being slowly destroyed in winter. They are being destroyed at what for them
is their optimum and maximum rate.

\section{First Order and Zero Order Rates of Germination}

One of the principles that came out of physical organic chemistry in the 1940-
1960 era was that complex chemical processes involving a number of individual
chemical steps could still follow simple rate laws. Generally one step in the sequence
will be significantly slower than the rest. This step is called the rate determining (r. d.)
step. It not only limits the rate of the overall process, but the rate law of this rate
determining step becomes the rate law for the whole process. This discovery does not
seem to be appreciated by horticulturists and biologists. The old idea still prevails that
biological processes are so complex that they are not amenable to precise
mathematical treatment.

In the present work all data were treated as rate data. When sufficient data
were available, rate plots were constructed. The number or percent germination was
on the vertical coordinate and the time in days on the horizontal coordinate. Generally
these plots approached one of the four types illustrated by Figures 4-1 to 4-4 in
Chapter 4 for at least 60-80\% of the germination in any single cycle. These plots can
be expressed more concisely in tabular form. Tables 19-1 through 19-6 illustrate how
exact these fits can be to simple zero order and first order kinetics. Tables 19-1, 19-2,
and 19-3 show a fit with zero order kinetics, dn/dt = k. Tables 19-4, 19-5, and 19-6
show a fit with first order kinetics, dn/dt = kn. The closeness of the fit is shown by
closeness of the values in columns 2 and 3.

The examples in these Tables were chosen to show a wide range of conditions.
Table 19-1, 19-3, and 19-4 are for germinations taking place at 70. Table 19-2 is for a
germination taking place at 40. Table 19-5 is for a germination that takes place only in
light using alternating cycles of twelve hours dark and twelve hours light. Table 19-6 is
for a germination that requires gibberelic acid-3 (GA-3).
The good fits in Tables 19-1 through 19-6 could be dismissed as further
examples of a rate-determining step dictating a simple rate law in a complex process.
These are well known in mechanistic chemistry, but there we are dealing with
homogeneous systems which exchange energy continuously between the component
molecules. In contrast seeds are isolated systems and even their component cells
may be relatively isolated insofar as they exchange only small molecules such as
water and oxygen.

A limiting situation can be envisioned in which the seeds behave like
radioactive atomic nuclei which are similarly isolated. When the internal dynamics
achieve a certain configuration and distribution of energies, reaction takes place. This
would be germination with seeds and radioactive decomposition with atomic nuclei.

The rates of achievement of such reactive states are governed by Boltzmann statistics.
A certain percentage reacts in each time period generating first order kinetics.
To explain zero order kinetics for germination is more of a challenge. The
following is proposed. In a complex series of steps there are two slow steps. For
convenience let us regard them as consecutive and represent them as A to B to C.
Both are first order processes and both have about the same rate constant. It can be
shown mathematically that for the time period in which 20-90\% of A is being
destroyed, the amount of B remains roughly constant. During this period the formation
of B from A is roughly compensated by the conversion of B to C. This approximate
constancy of the amount of B means that C will be produced at a constant (zero order)
rate. Now if germination were to depend on the concentration of C or the
concentration of C reaching a critical level, the overall rate of germination would
approach zero order. This explanation is a bit too superficial. It will be for others to
provide more insight into the kinetics of germination, and in fact the kinetics of all
biological processes.

At the beginning of the work it was thought that there would be so much genetic
variation in seeds even from a single capsule that the rates of germination would
follow the laws of chemical kinetics in only an approximate fashion at best. In fact the
data fit zero order rates quite exactly in most experiments suggesting that most species
produce seeds that are genetically homogeneous, at least in respect to rates of
germination. In fact pure species generally produce seedlings that are
indistinguishable from each other, and the appearance of a mutation is relatively rare.

\begin{table}[h!]
\refstepcounter{table}
\begin{center}
Mimulus luteus - Table \thetable

Rate of Germination at 70, zero order rate,
ind. t 1.2 days, rate 17\% per day

\end{center}

\begin{tabular}{lll}
\hline
Time      & \multicolumn{2}{c}{\% Germ. (a)} \\
(days)    & Found & Calculated \\
\hline
1 & 0  & 0 \\
2 & 14 & 14\\
3 & 31 & 34\\
4 & 49 & 48\\
5 & 64 & 65\\
6 & 79 & 82\\
7 & 97 & 99\\
\hline
\end{tabular}
\end{table}

\begin{table}[h!]
\refstepcounter{table}
\begin{center}
Lilium pumilum - Table \thetable

Rate of Germination of at 40, zero order rate,
ind. t 33 days, rate 5.5\% per day
\end{center}

\begin{tabular}{lll}
\hline
Time      & \multicolumn{2}{c}{\% Germ. (a)} \\
(days)    & Found & Calculated \\
\hline
0-33 & 0 & 0 \\
34 & 8 & 8 \\
38 & 26 & 26 \\
42 & 50 & 47 \\
46 & 60 & 63 \\
50 & 87 & 87 \\
53 & 93 & 100 \\
\hline
\end{tabular}
\end{table}

\begin{table}[h!]
\refstepcounter{table}
\begin{center}
Pulsatilla vulgaris - Table \thetable

Rate of Germination at 70, first order rate,
md. t 13 days, half life 7 days.
\end{center}

\begin{tabular}{lll}
\hline
Time      & \multicolumn{2}{c}{\% Germ. (a)} \\
(days)    & Found & Calculated \\
\hline
0-13 & 0 & 0 \\
14 & 2 & 0 \\
16 & 32 & 33 \\
19 & 57 & 57 \\
24 & 68 & 70 \\
26 & 75 & 76 \\
38 & 92 & 91 \\
\hline
\end{tabular}
\end{table}

\begin{table}[h!]
\refstepcounter{table}
\begin{center}
Primula cortusoides - Table \thetable

Rate of Germination of at 70 first order
rate, md. t 10 days, half life 5 days.
\end{center}

\begin{tabular}{lll}
\hline
Time      & \multicolumn{2}{c}{\% Germ. (a)} \\
(days)    & Found & Calculated \\
\hline
0-10 & 0 & 0 \\
11 & 10 & 13 \\
13 & 40 & 33 \\
18 & 66 & 65 \\
27 & 87 & 89 \\
32 & 93 & 94 \\
63 & 99 & 100 \\
\hline
\end{tabular}
\end{table}

\begin{table}[h!]
\refstepcounter{table}
\begin{center}
Aruncus sylvestris - Table \thetable

Rate of Germination in light at 70,
zero order rate, md. t 20 days,
rate 7.5\% per
\end{center}

\begin{tabular}{lll}
\hline
Time      & \multicolumn{2}{c}{\% Germ. (a)} \\
(days)    & Found & Calculated \\
\hline
0-20 & 0 &  0 \\
22 & 17 & 15 \\
24 & 30 & 30 \\
26 & 46 & 45 \\
28 & 63 & 60 \\
30 & 76 & 75 \\
32 & 89 & 90 \\
\hline
\end{tabular}
\end{table}

\begin{table}[h!]
\refstepcounter{table}
\begin{center}
Echinocereus pectinatus - Table \thetable

Rate of Germination of with GA-3 at
70, first order rate, md. t 6 days,
half life 4.9 days
\end{center}

\begin{tabular}{lll}
\hline
Time      & \multicolumn{2}{c}{\% Germ. (a)} \\
(days)    & Found (a) & Calculated \\
\hline
0-5 & 0 & 0 \\
9.9 & 52 & 50 \\
14.8 & 76 & 75 \\
19.7 & 87.5 & 87.5 \\
24.6 & 93.8 & 93.8 \\
29.5 & 97.0 & 96.9 \\
34.4 & 98.3 & 98.4 \\
\hline
\end{tabular}

\hfil

(a) The precision expressed is real as a sample of 1200 seeds was used.
\end{table}

Rate data can be used to detect situations where two kinds of seeds are
present. An example was in the germination at 40 of a large sample of seed from a
colony of hybrids of Primulas of the Vernales section. The rate plot was a composite of
two zero order rate plots. One had a rate of 6\%/d and and ind. t of 16 d. The other was
much slower with a rate of 1 0/cld but with the same ind. t of 16 d.

\section{Induction Periods and Chemical Time Clocks}

Most of the rate plots were characterized by an induction period. The induction
times ranged generally from five days to fifty-five days. As discussed in Chapter 1 and
3, a significant induction period followed by an abrupt onset of germination is typical of
a chemical time clock and the destruction of a chemical inhibitor. They cannot be due
to time required for diffusion of water and/or oxygen since these would produce a
gradual onset of germination. From the viewpoint of survival, time clocks are another
mechanism for delaying germination until the seed is dispersed.

A word of explanation about chemical clock reactions is appropriate. These are
not common, but are familiar to most chemistry students as they are usually
demonstrated in the first course in college chemistry. The one most often chosen for
demonstration involves the near instantaneous change of a colorless solution to a
dark brown solution because of the formation of iodine. The induction time (delay
period) required for this change can be varied from seconds to minutes by varying the
concentrations of the reagents. Another and more common example of a chemical
time clock is present in all the articles made of rubber that surround us everywhere as
described in Chapter 4. All clock reaction involve the destruction of a critical block or
inhibitor. When this block is finally reduced below a critical level, the chemical process
starts up and quickly runs to completion. The chemical nature of these clock reaction
in plants have not been studied.

\section{Low Temperature Metabolism and Large and Negative Temperature
Coefficients for Rates of Metabolic Processes}

It has long been known that warm blooded animals control their temperature
and operate at essentially a constant temperature. When body temperature moves
above or below this temperature, rates of body reactions undergo large changes. It is
also known that arctic marine animals require low temperatures to live and that these
low temperatures must be within narrow limits. What has not been recognized is that
the chemical processes in plants often have equally enormous changes in rate with
temperature.

It has been found in the present work that where a cold cycle or a warm cycle is
required to destroy germination inhibitors, such processes take place either around 40
or around 70, but never at both temperatures. This means.that the change in rates
between 40 and 70 are very large, essentially infinite. .Thesurvivavalue of this is
apparent. Plants have evolved a winter metabolism and summer metabolism, and
generally these are kept separate and distinct. Nature plays no favorites here, and
there are about as many species in which the germination inhibitors are destroyed at
70 as at 40. Thus about as many undergo an enormous increase in rate on raising
temperature (positive temperature coefficient) as those that undergo an enormous
increase in rate on lowering temperature (negative temperature coefficient).
The above situation in biological systems is markedly different from the
situation in small molecule chemistry and in the usual chemistry laboratory. In small
molecule chemistry, rates of chemical reactions typically change by the order of 2-3 for
each ten degree change in temperature. Further, although reactions with negative
temperature coefficients are known, they are rare and arise in a different way. The
explanation of negative temperature coefficients in small molecule chemistry is shown
in the following sequence of steps 1 to 3.

\begin{itemize}
\item A+B = Intermediate
\item Intermediate = Inactive Species
\item Reactant + Intermediate = Product(s)
\end{itemize}

All three steps have positive temperature coefficients which is true of all \textit{single
step} chemical processes. It might seem impossible to combine three steps with
positive temperature coefficients anctget an overall process with a negative
temperature coefficient, but watch closely. What if step 2 slows down with falling
temperature much more than steps 1 and 3, and what if this factor is so large as to
dwarf the others. Now the concentration of the \textit{Intermediate} increases and builds up.
This causes step 3 to go faster leading to an overall increase in rate for the overall
combination of steps 1-3.

The fact that enormous temperature coefficient and negative temperature
coefficients, which are so rare in small molecule reactions, become so common in
biological processes means that biological systems have some special way to
generate these situations. The explanation almost certainly involves changes in the
configurations of enzymatic proteins. Many examples are known where proteins
change conformation over a very narrow range of temperatures (refs. 41 and 42), and
such changes can be reversible. If on lowering the temperatUre a critical. protein
changes its conformation to form either an enzymatically active form or a form that
allows some critical diffusion, this would explain the large temperature coefficients,
and one that can be either positive or negative.

Although chemical processes can exhibit a negative temperature coefficient
over a limited range of temperatures, this cannot persist indefinitely. Ultimately a
temperature is reached where further decreases in temperature slows down the
process. Thus low temperature metabolism will always exhibit an optimum
temperature in a plot of rate versus temperature This has been demonstrated in arctic
animals. It has now been demonstrated in the inhibitor destruction reactions in seed
germination. Although further work is needed to settle the matter conclusively, it is
suspected that somewhere between 40 and 70 the rates of inhibitor destruction
reactions undergo an enormous change in rate over a few degrees of temperature.
The rates of germination may go faster at 40 than 70 or the reverse. In this
respect they are like inhibitor destruction reactions. However, the change in rate
between 40 and 70 may or may not be large. Examples were found where the rates of
germination were identical at 40 and 70 or differed by only small amounts. In other
examples there was a large difference in rate of germination between 40 and 70.

\section{Optimum Rates Under Conditions of Oscillating Temperatures}

In Chapter 10 data was presented showing that some species give significant
germination only under outdoor conditions. The significant feature of these outdoor
conditions that set them apart from indoor treatments was the oscillating temperatures
that occurred each day. Although I know of no prior reports where a chemical process
was recognized to maximize its rate under conditions.of oscillating temperatures, the
phenomenon is theoretically possible. Suppose that in a complex series of reactions
there is a critical step that takes place only at 70 and another critical step that takes
place only at 40. In order for the overall process to take place the system must oscillate
between 70 and 40. Part of the time is spent at 70 to let the first process function and
part of the time must be spent at 40 to let the second process function. Only in this way
can the overall process take place.

We can only surmise at this time that biological processes offer more
opportunity for the above phenomenon, because of the prevalence of reactions that
have enormous changes in rates over narrow ranges of temperature.

